\section{Câu 3}
Cho mạch khuếch đại tín hiệu như hình vẽ. Các tụ $C_1$, $C_2$ có giá trị rất lớn. BJT có hệ số $\beta = 100$ và có mã là 2N2907

\begin{figure}[h!]
    \centering
    \resizebox{0.7\textwidth}{!}{
    \begin{tikzpicture}[transform shape, american voltages]
	% Paths, nodes and wires:
	\draw (8, 6) to[american resistor, l={$R_S=1k \Omega$}] (9.2, 6);
	\draw (10, 6) to[capacitor, l={$C_{C1}$}] (11, 6);
	\draw (12, 5.2) to[american resistor, mirror, l_={$R_2=10k \Omega$}] (12, 3.8);
	\draw (16, 4.6) to[capacitor, mirror, l_={$C_{C2}$}] (15, 4.6);
	\draw (7, 4) to[sinusoidal voltage source, l={$v_s$}] (7, 3);
	\draw (7, 4) -| (7, 6) -- (8, 6);
	\draw (9.2, 6) -- (10, 6);
	\draw (12, 8.2) to[american resistor, mirror, l_={$R_1 =10k \Omega$}] (12, 6.8);
	\draw (14, 4.2) to[american resistor, mirror, l={$R_E = 500\Omega$}] (14, 2.8);
	\node[ocirc](N1) at (18, 4.6){} node[anchor=west] at (N1.east){$v_o$};
	\node[npn] at (14, 6){};
	\draw (11, 6) -- (13.16, 6);
	\draw (12, 6.8) -| (12, 6);
	\draw (12, 6) -| (12, 5.2);
	\draw (14, 5.23) -| (14, 4.2);
	\draw (17, 4.2) to[american resistor, mirror, l={$R_L = 300\Omega$}] (17, 2.8);
	\draw (12, 8.2) -| (12, 9) -- (14, 9) -| (14, 6.8);
	\draw (12, 3.8) |- (14, 2) -| (14, 2.8);
	\draw (14, 4.6) -- (15, 4.6);
	\draw (16, 4.6) -- (18, 4.6);
	\draw (17, 4.6) -| (17, 4.2);
	\node[ground] at (17, 2){};
	\node[ground] at (7, 2){};
	\draw (14, 9.6) -| (14, 9);
	\draw (14, 2) -| (14, 1.4);
	\draw (17, 2) -| (17, 2.8);
	\draw (7, 3) -| (7, 2);
	\node[ocirc](N2) at (14, 9.6){} node[anchor=south] at (N2.north){$V^+ = +9V$};
	\node[ocirc](N3) at (14, 1.4){} node[anchor=north] at (N3.south){$V^- = -9V$};
	\node[circ] at (12, 6){};
	\node[circ] at (14, 2){};
	\node[circ] at (14, 9){};
	\node[circ] at (14, 4.6){};
	\node[circ] at (17, 4.6){};
\end{tikzpicture}%
    }
\end{figure}
\textbf{a) Vẽ VTC của mạch (kiểm chứng sử dụng mô phỏng) và tìm điểm hoạt động $Q$ của BJT.}

Theo DC mạch có dạng:
\begin{figure}[htbp]
\centering
\begin{subfigure}{0.48\textwidth}
  \centering
  \begin{tikzpicture}[transform shape]
	% Paths, nodes and wires:
	\draw (7, 4.198) to[american resistor, mirror, l_={$R_2=10k \Omega$}] (7, 2.798);
	\draw (7, 7.198) to[american resistor, mirror, l_={$R_1 =10k \Omega$}] (7, 5.799);
	\draw (9, 3.198) to[american resistor, mirror, l={$R_E = 500\Omega$}] (9, 1.799);
	\node[npn] at (9, 4.999){};
	\draw (7, 5.799) -| (7, 4.999);
	\draw (7, 4.999) -| (7, 4.198);
	\draw (9, 4.228) -| (9, 3.198);
	\draw (7, 7.198) -| (7, 7.998) -- (9, 7.998) -| (9, 5.799);
	\draw (7, 2.798) |- (9, 0.999) -| (9, 1.799);
	\draw (9, 8.599) -| (9, 7.998);
	\draw (9, 0.999) -| (9, 0.398);
	\node[ocirc](N1) at (9, 8.599){} node[anchor=south] at (N1.north){$V^+ = +9V$};
	\node[ocirc](N2) at (9, 0.398){} node[anchor=north] at (N2.south){$V^- = -9V$};
	\node[circ] at (7, 4.999){};
	\node[circ] at (9, 0.999){};
	\node[circ] at (9, 7.998){};
	\node[circ] at (9, 3.75){};
	\draw (7, 4.999) -- (8.16, 4.999);
	\node[ocirc](N3) at (11, 3.75){} node[anchor=west] at (N3.east){$v_o$};
	\draw (9, 3.75) -- (11, 3.75);
\end{tikzpicture}
%   \includegraphics[width=\linewidth]{fig1.png}
    \caption{Mạch ở chế dộ DC}
%   \label{fig:first}
\end{subfigure}
\hfill
\begin{subfigure}{0.48\textwidth}
  \centering
  \begin{tikzpicture}[transform shape]
	% Paths, nodes and wires:
	\draw (7.15, 5) to[american resistor, mirror, l_={$R_{th}=5k \Omega$}] (5.75, 5);
	\draw (9, 3.198) to[american resistor, mirror, l={$R_E = 500\Omega$}] (9, 1.799);
	\node[npn] at (9, 4.999){};
	\draw (9, 4.228) -| (9, 3.198);
	\draw (9, 6.5) -| (9, 5.769);
	\draw (9, 1.799) -| (9, 1.198);
	\node[ocirc](N1) at (9, 6.5){} node[anchor=south] at (N1.north){$V^+ = +9V$};
	\node[ocirc](N2) at (9, 1.198){} node[anchor=north] at (N2.south){$V^- = -9V$};
	\node[circ] at (9, 3.75){};
	\node[ocirc](N3) at (11, 3.75){} node[anchor=west] at (N3.east){$v_o$};
	\draw (9, 3.75) -- (11, 3.75);
	\node[ground] at (5, 4){};
	\draw (5, 4) -| (5, 5) -- (5.75, 5);
	\draw (7.15, 5) -| (8.16, 4.999);
\end{tikzpicture}
%   \includegraphics[width=\linewidth]{fig2.png}
    \caption{Mạch sau khi đã Thevenin}
%   \label{fig:second}
\end{subfigure}
% \caption{Overall caption for both}
% \label{fig:both}
\end{figure}

Thevenin ta có:
\[ R_{th} = R_{1} // R_{2} = 5\,\textsf{k}\Omega \]
\[ V_{th} = \dfrac{R_{2}}{R_{1} + R_{2}} V_{+} + \dfrac{R_{1}}{R_{1} + R_{2}}V_{-} = 0\,\textsf{V} \]

Áp dụng KVL cho vòng từ cực B tới cực E:
\[ I_{B}R_{th} + V_{BE} + I_{E}R_{E} - 9 = 0 \]

Ta có: $I_E = (\beta + 1)I_B$
\[
\Longrightarrow I_B = \frac{9 - V_{BE}}{R_{th} + (\beta + 1)R_E}
= \frac{9 - 0.7}{5 + (100 + 1)\times 0.5}
= 0.1495\,\text{mA}
\]

\[
\Longrightarrow I_C = \beta I_B = 100 \times 0.1495\ \text{mA}
= 14.95\,\text{mA}
\]


Áp dụng KVL cho vòng từ cực C tới cực E:
\[-9+V_{CE}+I_E R_E-9=0 \]
\[\Rightarrow V_{CE}=18-(\beta+1)I_B\times R_E=10.4477\ V\]

Vậy điểm làm việc $Q$ là: $Q(I_{CQ},V_{ECQ}) = Q(14.95\ mA,\ 10.4477\ V)$

\textbf{b) Đặt $v_{s} = V_{m} \sin \left( \omega t\right)$ vào mạch. Tìm $A_{vo}$, $G_v$, $R_i$, $R_o$ của mạch.}

Mạch ở chế độ AC có dạng:
\begin{figure}[h!]
    \centering
    \resizebox{0.7\textwidth}{!}{
    \begin{tikzpicture}[transform shape]
	% Paths, nodes and wires:
	\draw (8, 3.25) to[american resistor, mirror, l_={$R_{th}=5k \Omega$}] (8, 4.65);
	\draw (11, 3.25) to[american resistor, mirror, l_={$R_E = 500\Omega$}] (11, 1.85);
	\node[npn] at (11, 4.999){};
	\node[circ] at (11, 3.75){};
	\node[ocirc](N1) at (16, 3.75){} node[anchor=west] at (N1.east){$v_o$};
	\draw (11, 3.75) -| (16, 3.75);
	\node[ground] at (8, 1){};
	\draw (5, 3.25) to[american resistor, mirror, l_={$R_{S}=1k \Omega$}] (5, 4.65);
	\node[oscillator](N2) at (3.49, 3.99){} node[anchor=west] at (N2.east){$v_s$};
	\draw (14, 3.25) to[american resistor, mirror, l_={$R_L = 300\Omega$}] (14, 1.85);
	\node[ground] at (3, 1){};
	\node[ground] at (5, 1){};
	\node[ground] at (11, 1){};
	\node[ground] at (14, 1){};
	\node[ground] at (13, 5.75){};
	\draw (3, 3.5) -| (3, 1);
	\draw (3, 4.48) |- (3, 5) -| (10.16, 4.999);
	\draw (5, 5) -| (5, 4.65);
	\draw (5, 3.25) -| (5, 1);
	\draw (8, 1) -| (8, 3.25);
	\draw (8, 4.65) |- (8, 5);
	\draw (11, 5.769) -| (11, 6.25) -- (13, 6.25) -| (13, 5.75);
	\draw (11, 4.229) -| (11, 3.25);
	\draw (11, 1.85) -| (11, 1);
	\draw (14, 1) -| (14, 1.85);
	\draw (14, 3.25) -| (14, 3.75);
	\node[circ] at (5, 5){};
	\node[circ] at (8, 5){};
	\node[circ] at (14, 3.75){};
\end{tikzpicture}
}
\end{figure}

Xét hoạt động chế độ AC cho toàn mạch:

\[
r_\pi = \frac{V_T}{I_B} = \frac{\beta V_T}{I_C} = 0.169\ k\Omega
\]

\[
r_o = \frac{V_A}{I_C} = \frac{100}{14.8} = 6.76\ k\Omega
\]

\[
A_{vo} = \frac{R_E}{\dfrac{r_\pi}{\beta + 1} + R_E} 
= \frac{0.5}{\dfrac{0.169}{101} + 0.5} 
= 0.997\ V/V
\]

\[
R_{in} = R_1 // R_2 // (\beta + 1)\left(r_e + R_E // r_o\right) = 10k // 10k // 101 \times \left(\frac{r_\pi}{\beta + 1} + 500 // 6.76k\right)
= 4.52\ k\Omega
\]

\[
R_o = R_E // r_e // r_o 
= 0.5k // \frac{0.167k}{101} // 6.76k 
= 1.65\ \Omega
\]

\[
G_v = \frac{R_{in}}{R_{in} + R_{sig}} \times 
\frac{R_L}{R_L + R_o} \times A_{vo} 
= \frac{4.52}{4.52 + 1} \times 
\frac{300}{300 + 1.65} \times 0.997 
= 0.81\ V/V
\]



\textbf{c) Lựa chọn các tụ $C_{1}$, $C_{2}$ để mạch có $f_{L}=100\,\textsf{Hz}$.}